\documentclass[11pt]{article}

\usepackage[T1]{fontenc}
\usepackage[utf8]{inputenc}
\usepackage{lmodern}
\usepackage{microtype}
\usepackage[a4paper,margin=30mm]{geometry}
\usepackage{amsmath,amssymb,amsthm,mathtools}
\usepackage[hidelinks]{hyperref}
\usepackage[nameinlink,noabbrev]{cleveref}

\newtheorem{theorem}{Theorem}
\newtheorem{corollary}[theorem]{Corollary}
\newtheorem{lemma}[theorem]{Lemma}
\theoremstyle{definition}
\newtheorem{definition}[theorem]{Definition}
\theoremstyle{remark}
\newtheorem{remark}[theorem]{Remark}

\newcommand{\N}{\mathbb{N}}
\newcommand{\Z}{\mathbb{Z}}
\newcommand{\FF}{\mathbb{F}}
\newcommand{\card}[1]{\lvert #1\rvert}
\newcommand{\Tri}{\operatorname{Tri}}
\newcommand{\dens}{\operatorname{d}}

\title{Sparse Deletion and Triangular Prefix Sums}
\author{Przemys\l aw Czerkas}
\date{\today}

\begin{document}

\maketitle

\begin{abstract}
Let $E\subseteq\N$ have natural density zero, and let
$a_0<a_1<a_2<\cdots$ enumerate the positive integers outside $E$.
We prove that the set of indices $n$ for which
\[
  \sum_{i=0}^{n-1} a_i^2
\]
is a triangular number has natural density zero.  The result permits an
arbitrary density-zero deletion set; no regularity hypothesis is imposed on
its elements or their locations.  The proof combines a finite-field sieve,
the Chinese remainder theorem, a sparse-perturbation principle, and a transfer
from value indices to indices in the retained enumeration.  A Lean~4/Mathlib
formalization of the theorem, including an axiom audit with no use of
\texttt{sorryAx}, is available at
\url{https://github.com/pczerkas/sparse-deletion-lean}.
\end{abstract}

\noindent\textbf{Keywords.} natural density; triangular numbers; sums of
squares; finite fields; Chinese remainder theorem; formal verification; Lean.

\noindent\textbf{MSC 2020.} 11B75; 11T06; 11N35; 03B35.

\section{Introduction}

A nonnegative integer $s$ is \emph{triangular} if
\[
  2s=m(m+1)
\]
for some $m\in\N$.  Write $\Tri(s)$ for this condition.  For a set
$A\subseteq\N$, let
\[
  A(x)=\card{A\cap\{0,1,\ldots,x-1\}}.
\]
We say that $A$ has natural density zero if $A(x)/x\to0$ as $x\to\infty$.

The purpose of this note is to establish a stability result for triangular
prefix sums of squares under arbitrary sparse deletion.  The point is that a
density-zero deletion can have very large gaps and can delete very large
integers; the argument does not require a quantitative bound on either
phenomenon.

\begin{theorem}[Sparse deletion theorem]\label{thm:main}
Let $E\subseteq\N$ have natural density zero, and enumerate the positive
integers outside $E$ increasingly as
\[
  a_0<a_1<a_2<\cdots.
\]
Then
\[
  \dens\biggl\{n\in\N:
    \Tri\!\left(\sum_{i=0}^{n-1}a_i^2\right)\biggr\}=0.
\]
\end{theorem}

For $n=0$, the sum is interpreted as $0$, which is triangular.  This isolated
index clearly does not affect the conclusion.

\begin{remark}
\Cref{thm:main} is a density statement, not a finiteness statement: it neither
asserts that only finitely many triangular prefix sums occur nor that they
eventually disappear.  It also does not apply directly to prefix sums over
primes, since the complement of the primes is not a density-zero deletion
set.
\end{remark}

\section{Proof architecture}

We record the ingredients of the proof, both to explain the theorem and to
make clear the division between its arithmetic and density-theoretic parts.
The complete formal proof is in the accompanying Lean repository.

\subsection{A cubic modular obstruction}

Put
\[
  S(x)=\sum_{k=0}^{x}k^2
      =\frac{x(x+1)(2x+1)}{6}.
\]
After an affine change of variables over a finite field of characteristic
neither $2$ nor $3$, the condition that $S(x)+b$ be triangular is controlled
by whether a value of a cubic polynomial is a square.  The relevant cubic is
affinely equivalent to
\[
  t\longmapsto t^3-t.
\]

The required value-set estimate is elementary.

\begin{lemma}[Cubic value-set bound]\label{lem:cubic}
Let $\FF_q$ be a finite field of characteristic different from $3$.  Then
\[
  \card{\{t^3-t:t\in\FF_q\}}\geq \frac{2q-1}{3}.
\]
\end{lemma}

In the formal development, \Cref{lem:cubic} follows by studying collisions
of $t^3-t$.  Distinct elements in a common fibre give points on the conic
\[
  x^2+xy+y^2=1,
\]
and an elementary injection bounds the number of such collisions.  No Weil
bound, Hasse bound, elliptic-curve estimate, or character-sum estimate is
used.

Together with the elementary count of quadratic residues in $\FF_p$, this
implies the following uniform local sieve bound.  There is an infinite sequence
of pairwise distinct primes $p\geq37$ such that, for every integer shift $b$,
at most a proportion $6/7$ of residue classes modulo $p$ can satisfy the
necessary modular condition for $S(x)+b$ to be triangular.  The same estimate
holds after translating $x$ by any integer.

\subsection{CRT amplification}

Choose $r$ distinct primes from the preceding family.  Since their moduli are
pairwise coprime, the Chinese remainder theorem multiplies the local bounds.
Thus, for every integer shift $b$, there is a positive block length $M$ such
that every interval $[u,u+M)$ contains at most
\[
  \left(\frac67\right)^rM
\]
integers $x$ for which $S(x)+b$ can be triangular.  Crucially, $M$ works
uniformly for the starting point $u$ and for the shift $b$.

\subsection{Sparse perturbations}

For the deletion set $E$, define
\[
  D(x)=-\sum_{\substack{0\leq k\leq x\\ k\in E}}k^2.
\]
Then the sum of squares of retained positive integers up to $x$ is represented
by $S(x)+D(x)$, up to the harmless indexing convention at $0$.  The function
$D$ is locally constant away from $E$: its increment can occur only at a
deleted integer.  The following abstract statement supplies the density step.

\begin{lemma}[Sparse-perturbation principle]\label{lem:perturbation}
Let $E\subseteq\N$ have natural density zero.  Let $f,d:\N\to\Z$ and
$A\subseteq\Z$, and suppose that $d(n+1)=d(n)$ whenever $n+1\notin E$.
Assume that, for every $\varepsilon>0$, there are $M>0$ and $C\in\N$ with
$C/M<\varepsilon$ such that, for all $u\in\N$ and $b\in\Z$,
\[
  \card{\{n\in [u,u+M):f(n)+b\in A\}}\leq C.
\]
Then
\[
  \dens\{n\in\N:f(n)+d(n)\in A\}=0.
\]
\end{lemma}

Apply \Cref{lem:perturbation} with $f=S$, with the perturbation $D$, and
with $A$ the set of triangular integers.  The CRT estimate makes $C/M$
arbitrarily small by increasing $r$, hence triangular hits have density zero
when indexed by their original values.

Finally, a density-zero deletion leaves sufficiently many integers: the
$n$-th retained positive integer is eventually less than $2n+2$.  An injection
from triangular hits in the retained enumeration to the corresponding hits at
value indices transfers density zero to the index set in \Cref{thm:main}.
This completes the proof.

\section{A concrete corollary}

Let $c_0<c_1<c_2<\cdots$ enumerate the composite positive integers.  Since
the set consisting of $1$ and all primes has natural density zero,
\Cref{thm:main} yields the following consequence.

\begin{corollary}
The set
\[
  \left\{n\in\N:
    \Tri\!\left(\sum_{i=0}^{n-1}c_i^2\right)\right\}
\]
has natural density zero.
\end{corollary}

\section{Formal verification and reproducibility}

The formalization is written in Lean~4 \cite{lean4} and uses Mathlib
\cite{mathlib}.  The repository pins
Lean and Mathlib to version \texttt{v4.35.0}; the exact dependency revisions are
recorded in \texttt{lake-manifest.json}.  The main formal theorem is
\begin{verbatim}
SparseDeletion.main_sparse_deletion_unconditional
    (E : Set Nat) (hE : HasDensityZero E) :
    HasDensityZero {n : Nat | Triangular (keptIndexPrefix E n)}
\end{verbatim}

The repository provides separate verification targets for the end-to-end
theorem and the generic sparse-perturbation lemma.  It also contains targets
for the exact finite-field value-set formula and the composite-number
corollary.  With Lean/Lake installed, run:
\begin{verbatim}
git clone https://github.com/pczerkas/sparse-deletion-lean.git
cd sparse-deletion-lean
lake exe cache get
lake build
\end{verbatim}
The source includes axiom-printing certificate files.  The principal proof
terms use only the standard Lean axioms \texttt{propext},
\texttt{Classical.choice}, and \texttt{Quot.sound}; they do not depend on
\texttt{sorryAx}.  As usual, formal verification establishes correctness
relative to the Lean kernel, the pinned dependency code, and the accuracy of
the formalized statement.

The repository also formalizes an optional, stronger classical finite-field
identity: for $q\equiv1\pmod3$ (respectively $q\equiv2\pmod3$), the image of
$t\mapsto t^3-t$ in $\FF_q$ has cardinality $(2q+1)/3$ (respectively
$(2q-1)/3$).  This exact formula is not needed for \Cref{thm:main}.

The source bundle includes \texttt{figure-proof-architecture.txt}, a
plain-text dependency sketch suitable for conversion into an editorial figure
if desired.

\section*{References}
\begin{thebibliography}{9}

\bibitem{lean4}
Leonardo de Moura and Sebastian Ullrich.
\newblock The Lean 4 theorem prover and programming language.
\newblock In \emph{Automated Deduction -- CADE 28}, pages 625--635,
Springer, 2021.
\newblock \url{https://doi.org/10.1007/978-3-030-79876-5_37}.

\bibitem{mathlib}
The mathlib Community.
\newblock The Lean Mathematical Library.
\newblock In \emph{Proceedings of the 9th ACM SIGPLAN International
Conference on Certified Programs and Proofs (CPP 2020)}, pages 367--381,
ACM, 2020.
\newblock \url{https://doi.org/10.1145/3372885.3373824}.

\end{thebibliography}

\section*{Availability and licence}

The complete Lean source, build instructions, and dependency audit are publicly
available at
\begin{center}
\url{https://github.com/pczerkas/sparse-deletion-lean}.
\end{center}

This manuscript is licensed under the
\href{https://creativecommons.org/licenses/by/4.0/}{Creative Commons
Attribution 4.0 International (CC BY 4.0) licence}.  The accompanying Lean
software repository is distributed separately under its repository licence
(MIT at the time of this manuscript).

\section*{AI-use declaration}

AI systems, including OpenAI ChatGPT, were used for mathematical exploration,
proof development, refactoring, and Lean assistance.  The author reviewed the
mathematical statement and the formal development.  The machine-checkable
formal proof and its pinned dependencies are available in the repository cited
above.

\end{document}
